\scanpage{188}
% archive sheet number: 187

$A$ algèbre affine d'un groupe affine $G$ (pas néc. de type fini) sur l'anneau $k$

\[G\unsure{\longrightarrow|\text{---}}\unsure{G_m|\mathbb G_m}\quad\text{hom de $k$-grp}\quad\unsure{\approx|\simeq}\quad G\longrightarrow\mathbb G_a\]
morphisme de $k$-schémas, tel que
\begin{enumerate}[label=(\roman*)]
\item Addition transformée en multiplication
\item $0$ s'en va en $1$
\end{enumerate}

i.e. $f\in A$, telle que
\[\Delta f=f\otimes f\tag*{a)}\]
\[\varepsilon(f\unsure{)|\}}=1\tag*{b)}\]

Supposant que $A$ soit \uline{libre}, alors $A$ est l'ensemble des formes $k$-linéaires continues sur $\check A=\operatorname{Hom}(A,k)$, muni de la convergence simple. Alors a) signifie que $f$ est multiplicatif \deleted{\unsure{et}} [que $\check A$ est muni de la structure de $k$-algèbre déduite de $\Delta$] et b) que $f$ est unitaire; donc \unsure{a/. et b|c/. et b}, que $f$ est un hom d'algèbres topologiques.

\deleted{D'où, si} Note que rien ne garantit que $\check A$ soit topologisé linéairement en tant qu'anneau, i.e. qu'il y ait un système fondamental de voisinages qui soient des \uline{idéaux}; \unsure{or|ou}, ce qui revient au $=$ si $k$ est un corps [et ce qui suffit en tous cas]. rien ne
